% a06_BALtorus_d7: part of Step 3 of the proof of thm:a06:T1 (input from sec/a06_BALtorus_c.tex).
% Group (D): Ba3 Prop. 1 at L = 5, second half: the renormalisation of the divergent classes (pp. 434-445) and the assembled statement; the one open row.

\subsubsection*{Step 3(D$_4$): the renormalisation of the divergent classes, \cite[pp.~434--445]{Ba3}}

After the cube steps of \cite[p.~433]{Ba3} a class component has scalar lines $G_k(0)$ on $\eta\Z^d$, sliced by \cite[(2.6)]{Ba3}, photon lines on the torus, vertices in the cube's inner part, all other objects external, and per-label full classes. Orderings are the lexicographic ones of the slice indices; for an ordering $\tilde l$ the components $G^{(\alpha)}_i$ along it form a laminar family. A \emph{dangerous node} is a component with $D\le0$; by the degree formula \cite[(2.1)--(2.3)]{Ba3} it is one of: a \cite[(2.4)]{Ba3} subgraph ($D=0$); $2\phi$ ($D\in\{-1,0\}$, $D=-1$ only at one loop, \cite[(3.6), (3.8)]{Ba3}); $2A$ ($D\in\{-1,0\}$, $D=-1$ only at one loop, \cite[(3.25)]{Ba3}); $\phi\phi A$ ($D=0$, one loop: $D\ge V/2+E/2-3$ with $E=3$ forces $V\le3$ and one loop); a vacuum component, cancelled by the vacuum counterterms \cite[(1.24)]{Ba3} with the same labels and lines, up to the torus-versus-$\eta\Z^d$ image terms. Subgraphs with no external scalar leg and an odd number of external vector legs vanish as factors at $B=0$ (their scalar index loops close inside them, $\Tr q^{\mathrm{odd}}=0$). A one-particle-reducible $2\phi$ node is not partnered (\cite[p.~430]{Ba3}: ``We consider one-particle-irreducible graphs only''); since $D(P_1+\text{line}+P_2)=D(P_1)+D(P_2)+2$, a dangerous reducible $2\phi$ chain consists of one-loop $D=-1$ parts, and after these are processed the chain is no longer dangerous. External vector fields inside vertices ($\tilde A$, and $\tilde B'$ after the cube step) count as external vector legs for the type of a node (an $\tilde A$-leg adds $\kappa=n'=1$ where an external leg adds $1$ to $E$), so the partnered $2\phi$ nodes are exactly the $2\phi$ nodes without such legs. A node $S$ is \emph{useful} along $\tilde l$ if its internal lines precede all its external lines, i.e.\ if it is a component. \emph{Processing} a useful dangerous node $S$ with an anchor $\hat x$ (a vertex of $S$): \cite[(2.4)]{Ba3} subgraphs are handled inside the smallest dangerous non-\cite[(2.4)]{Ba3} node containing them (\cref{lem:a06d:S44}), or alone by \cite[(2.8)]{Ba3} at either vertex; $2\phi$: the pairing of $G$ with its counterterm member, anchored at the vertex fixed by the counterterm convention (Agr), the exact identity \cite[(3.19)]{Ba3} for $D=0$ and \cite[(3.11)]{Ba3} for $D=-1$, which leaves a first-order local vertex with the regrouped coefficient and a remainder of degree $1+\alpha$; $2A$, $\phi\phi A$: Taylor expansion of the external fields at $\hat x$ (zeroth order for $D=0$, first order for $D=-1$), with \cref{lem:a06d:S27} for the multipliers and \cref{lem:a06d:S46} for sharp indicators, the local terms carrying the regrouped coefficients of \cref{thm:a06d:S48} and the remainders gaining $\alpha$ (resp.\ $1+\alpha$) on a path inside $S$. Non-\cite[(2.4)]{Ba3} inner nodes are processed before outer ones (\cite[p.~445]{Ba3}: ``assuming that it holds for all its subgraphs''), and the Taylor split of an outer node is applied to all terms produced by the inner processing (\cref{lem:a06d:S42}).

\begin{lemma}[exactness and degree preservation]\label{lem:a06d:S41}
\textup{(a)} Each processing step is an identity, at fixed external positions, of the sum over the orderings in which $S$ is useful: the Taylor formulas, \cite[(3.19), (3.26), (3.34), (2.8)]{Ba3} and the regrouping into classes are identities, and the counterterm pairing is an identity of class sums (\cref{thm:a06d:S23}). Hence the local coefficient of $S$, summed over $S$'s internal orderings, has the value of the regrouped \emph{unprocessed} coefficient; more precisely, the processed expression evaluated at external fields of $S$ that are polynomials of degree $\le1$ in the positions equals the unprocessed one, because every identity holds for every external field. \textup{(b)} Processing a useful node $S$ changes degrees only for subgraphs of $S$ containing a line of the Taylor/H\"older path ($+\alpha$, or $+1+\alpha$) and for subgraphs containing an external line of $S$ that carries the H\"older quotient ($-\alpha$, or $-1-\alpha$); every component along $\tilde l$ that contains such an external line contains $S$, so its net change is $0$. In the local term $S$ becomes a vertex $v_S$ at $\hat x$ with vertex power $D(S)$ at its vertex scale $s(v_S)=L^{j_0}\eta$ ($j_0$ the lowest external index), and every component containing $v_S$ has the degree of the original component with $S$, which is the convention \cite[(2.3)]{Ba3}.
\end{lemma}

\begin{proof}
(a) as listed; the inner processings are identities on complete groups of inner orderings, which are sub-orderings of $S$'s internal orderings. (b) A later component that contains a vertex of $S$ contains all of $S$, since $S$ is a component and components grow by merging along $\tilde l$; an earlier component contains no external line of $S$, these having higher index.
\end{proof}

\begin{lemma}[{inner processing acting on an outer node's external fields; \cite[p.~445]{Ba3}}]\label{lem:a06d:S42}
Let $S'$ be a dangerous non-\cite[(2.4)]{Ba3} node strictly inside a useful dangerous node $S$ along $\tilde l$, processed first with any anchor. In each term it produces, an external field $F$ of $S$ enters either unchanged or through a combination $\Delta[F]$ that vanishes for constant $F$ (a difference $F(x_i)-F(x_j)$, a Taylor remainder $F(x_i)-F(x_j)-(x_i-x_j)\cdot\partial F(x_j)$, or a derivative $\partial F(x_i)$); call the latter $\Delta$-terms. Then: \textup{(a)} $S$'s degree in a $\Delta$-term is $\ge D(S)+\alpha$: the gain of $S'$'s processing lies inside $S$ and the cost on $S$'s external line, which is not counted in $D(S)$; if $D(S)=0$, $S$ is not dangerous in a $\Delta$-term, and if $D(S)=-1$, $S$ is one-loop and contains no non-\cite[(2.4)]{Ba3} dangerous node. \textup{(b)} Apply $S$'s Taylor split (for $2\phi$, its counterterm pairing) to every term, $\Delta$-terms included: in a $\Delta$-term the local part is $\Delta$ evaluated at a constant field, i.e.\ $0$; for $2\phi$ the corresponding term of the counterterm member vanishes, since in $T_S(\cdot)$ the external fields of $S$ sit at $\hat x$ and $\Delta$ of them is $0$; the remainder of a $\Delta$-term is the $\Delta$-term itself. \textup{(c)} The local part of $S$, summed over all terms, is the processed expression evaluated at constant external fields of $S$, which by \cref{lem:a06d:S41}(a) equals the regrouped unprocessed coefficient; so \cref{thm:a06d:S48} bounds it, and the identities of \cref{lem:a06d:S26} and \cite[(3.30)]{Ba3} apply to it.
\end{lemma}

\begin{proof}
(a) is \cref{lem:a06d:S41}(b). (b) The external fields of $S$ are evaluated at $\hat x$ by the zeroth-order split, and every $\Delta$ annihilates constants; the counterterm member's statement is the definition of the counterterm graph \cite[p.~417]{Ba3}. (c) \cref{lem:a06d:S41}(a) at constant fields.
\end{proof}

So no anchor rule is needed: inner processings may act on outer fields; when they do, the outer node becomes convergent in those terms and its local coefficient is unchanged. Ba{\l}aban's choice (derivatives never on $G_0$'s legs) is one admissible choice among several. The anchors of $2A$ and $\phi\phi A$ nodes are free; the anchor of a $2\phi$ node is the vertex at which the counterterm graph of \cite[p.~417]{Ba3} merges the legs, the same in $G$ and in $T_SG$.

\begin{lemma}[{the \cite[(2.4)]{Ba3} subgraphs inside a node; \cite[p.~443]{Ba3}}]\label{lem:a06d:S44}
Let $S$ be a useful dangerous node of type $2\phi$, $2A$ or $\phi\phi A$ with anchor $\hat x$, and $\Sigma_1,\dots,\Sigma_r$ the fresh \cite[(2.4)]{Ba3} components inside $S$ along $\tilde l$ (vertex-disjoint, since a \cite[(1.8)]{Ba3} vertex has one derivative leg). There is an order of transport and a choice of H\"older paths such that in every term of the processed expression each $\Sigma_i$ either carries a Taylor/H\"older factor on its line, or has a vertex that carries no field external to $S$, at which \cite[(2.8)]{Ba3} is applied with no derivative on an external field of $S$.
\end{lemma}

\begin{proof}
Write $\Sigma_i=\{a_i,b_i\}$, with $a_i$ the anchor if the anchor is a vertex of $\Sigma_i$. Step~1: for each $i$ move every field external to $S$ that sits at $b_i$ onto $a_i$, to the Taylor order of $S$: $F(x_{b_i})=F(x_{a_i})[+(x_{b_i}-x_{a_i})\cdot\partial F(x_{a_i})]+$ remainder, the remainder carrying $|x_{b_i}-x_{a_i}|^\alpha$ (resp.\ $^{1+\alpha}$) on $\Sigma_i$'s line, which makes $\Sigma_i$'s degree $\ge\alpha>0$; the $\Sigma_i$ are vertex-disjoint, so these steps commute, and $\Delta$-fields are moved the same way. Step~2: transport all fields to $\hat x$ along any paths, as \cite[(3.19), (3.26), (3.34)]{Ba3} require (for $2\phi$ the two-step splitting $\phi_b(x')-\phi_b(x)=[\phi_b(x')-\phi_b(y)]+[\phi_b(y)-\phi_b(x)]$, $y$ the point of $a_i$). In every term produced by step~2, $b_i$ carries no field external to $S$: its fields were moved to $a_i$ in step~1, step~2 moves fields only from $a_i$ and other vertices, and fields at $\hat x$ are constants in $x_{b_i}$. \cite[(2.8)]{Ba3} at $b_i$ moves the derivative from $\Sigma_i$'s line onto $b_i$'s vector leg, onto $g$, or onto $b_i$'s other scalar leg, all internal to $S$ or carrying a constant field; the transformed component gets degree $1$ and later components through $b_i$ contain $\Sigma_i$'s line. If $\Sigma_i$ joins the two leg vertices of $S$ (the type \cite[(3.8)]{Ba3}, or the one-loop $2A$ bubble), step~1 is $S$'s own first-order Taylor expansion along $\Sigma_i$'s line. Inside a regrouped local coefficient there are no external fields at all (they are factored out at $\hat x$), so the integration by parts is free there.
\end{proof}

\begin{lemma}[{useless counterterms; \cite[p.~438]{Ba3}}]\label{lem:a06d:S45}
Let $\gamma$ be a $2\phi$ subgraph with external lines $l_a,l_b$, fix the indices of all lines outside $\gamma$, and let $j_*$ be the lexicographically smaller of the indices of $l_a,l_b$, $s_*\defeq L^{j_*}\eta$. Let $c_{\mathrm{use}}(x;j_*)$ be the counterterm coefficient of $T_\gamma$ restricted to the internal index assignments in which $\gamma$ is useless (maximal internal index $\ge j_*$), with $\gamma$'s inner divergences processed. Then $|c_{\mathrm{use}}(x;j_*)|\le Cs_*^{D(\gamma)}$ for $D(\gamma)=-1$ and $\le C(1+\log_L(1/s_*))\le C_{\alpha'}s_*^{-\alpha'}$ for $D(\gamma)=0$; consequently, in assignments where $\gamma$ is useless, $T_\gamma G$ is bounded by the bound function of $G/\gamma$ ($\gamma$ contracted to a vertex $v_\gamma$ at the anchor) with vertex power $D(\gamma)-\alpha'$ at $s(v_\gamma)=s_*$.
\end{lemma}

\begin{proof}
For a fixed top internal index $j_{\mathrm{top}}\ge j_*$ the bound function estimate of \cite[(2.16)]{Ba3}, with no sign condition on the degree and the anchor's cube fixed, bounds the internal sum by $Cs_{j_{\mathrm{top}}}^{D(\gamma)}$; it needs positive degrees for $\gamma$'s components along its orderings, which the inner processing supplies (induction on lines). Sum over $j_{\mathrm{top}}\in[j_*,k-1]$: for $D=-1$, $\sum_js_j^{-1}\le s_*^{-1}/(1-L^{-1})$; for $D=0$ the sum has $k-j_*=\log_L(1/s_*)$ terms. The lines of $G/\gamma$ at $v_\gamma$ are $l_a,l_b$, so $j(v_\gamma)=j_*$ and $s(v_\gamma)=s_*$, exactly the vertex scale of the bound function, and vertex powers that are powers of the vertex scale are allowed with any real exponent. In the useful assignments $T_\gamma G$ is paired with $G$ by \cite[(3.19)]{Ba3}/\cite[(3.11)]{Ba3} at the same indices (\cref{lem:a06d:S41}(a)); in the useless ones $G$ has no $\gamma$-divergence, since $\gamma$ is not a component. The loss $\alpha'\defeq\alpha/2$ is below the gains $\alpha$ of the H\"older remainders and below the half-integer gaps (degrees lie in $\frac12\Z$).
\end{proof}

\begin{lemma}[sharp indicators, boundary layers]\label{lem:a06d:S46}
Let $S$ be a dangerous node with $D(S)\le0$ whose local coefficient contains a sharp indicator $\mathbf 1_F(x_b)$ at an internal vertex $b$, $F$ the unit cube of a \cite[(1.6)]{Ba3}/\cite[(1.7)]{Ba3} vertex or $\Omega_1$. Write $c(\hat x)=\mathbf 1_F(\hat x)c^{\mathrm{unr}}(\hat x)+d(\hat x)$ with $c^{\mathrm{unr}}$ the coefficient with $\mathbf 1_F$ evaluated at $\hat x$. Then $|d(\hat x)|\le Cs_S^{D(S)}\exp(-\delta'\dist(\hat x,\partial F)/s_S)$, $s_S$ the top internal scale of $S$; summed over the cube $\Delta$ of side $s(v_S)$ of the local vertex, $\sum_{\hat x\in\Delta}\eta^d|d|\le Cs_S^{D(S)}(s_S/s(v_S))s(v_S)^d$; and after summing over $s_S\le s(v_S)$ the defect is bounded like a vertex of power $D(S)$ at $s(v_S)$, with no loss for $D(S)=0$ and better than \cite[(2.3)]{Ba3} for $D(S)=-1$.
\end{lemma}

\begin{proof}
$|\mathbf 1_F(x_b)-\mathbf 1_F(\hat x)|\le\mathbf 1\{\dist(\hat x,\partial F)\le|x_b-\hat x|\}$, the internal lines of $S$ at scales $\le s_S$ decay as $\exp(-\delta|x_b-\hat x|/s_S)$, and the bound function at fixed top scale gives the size; only the points of $\Delta$ within $O(s_S)$ of $\partial F$ contribute, a fraction $\le Cs_S/s(v_S)$ of $\Delta$; then a geometric sum.
\end{proof}

\begin{lemma}[position-space slices and remainder lines]\label{lem:a06d:S47}
\textup{(a)} Let $K$ be a kernel on $\xi\Z^3$ with $|\Delta_TD_SK(y,y')|\le C_K(\xi+\rho)^{-\omega-|S|-\alpha|T|}e^{-\delta\rho}$, $\rho\defeq|y-y'|$, for the derivative and H\"older operations used. With $\chi$ smooth, $=1$ on $[0,1]$, $0$ on $[L,\infty)$, $s_j=L^j\xi$, $\chi_0\defeq\chi(\rho/\xi)$, $\chi_j\defeq\chi(\rho/s_j)-\chi(\rho/s_{j-1})$ and the top piece $1-\sum_{j<J}\chi_j$ at $s_J=1$, the slices $K_j\defeq K\chi_j$ obey $|\Delta_TD_SK_j|\le C's_j^{-\omega-|S|-\alpha|T|}\exp(-\delta'\rho/s_j)$, so $K$ is a generalised line of the bound function with the derivatives acting on the sliced line. \textup{(b)} The remainder kernels $R=\breve G-C^\xi$ of the cube step ($|R|\le C$, $|\partial R|\le C(1+|\log(\xi+\rho)|)$, $|\partial R\partial^*|\le C(\xi+\rho)^{-1}$) have slices with bounds $s_j^{-(d-2)+1-|S|-\alpha|T|}$, one better than $C^\xi$-slices at every scale; for $|S|\ge1$ the logarithm is dominated by the terms where the derivative falls on $\chi_j$. \textup{(c)} For an inner node at scale $s_{j'}$ inside a coefficient on $\xi\Z^3$, the cumulative $C^\xi_{<j'}=C^\xi\chi_{<j'}$ splits as $C^{(j')}+R'_{j'}$ with $C^{(j')}\defeq(-\Delta^\xi+s_{j'}^{-2})^{-1}$, and after rescaling by $s_{j'}$, $R'_{j'}$ obeys the bounds of (b): $C^\xi-C^{(j')}=(s_{j'}^{-2}-1)C^\xi*C^{(j')}$ by the resolvent identity, and the cutoff part is supported at $\rho\sim s_{j'}$.
\end{lemma}

\begin{proof}
(a) On $\supp\chi_j$, $\rho$ is of order $s_j$ and $|\partial^k\chi_j|\le Cs_j^{-k}$; Leibniz. (b) Insert the bounds on $R$ into (a). (c) Resolvent identity and the convolution bound of two decaying kernels.
\end{proof}
